\subsection{Soft sheaves}

DEFINITION 5. --- Let \(p\colon E \rightarrow B\) be an étale map. One says that the map \(p\) is soft, or that the étale B-space \((E,p)\) is soft, if every continuous section of \(p\) over a closed subspace of B extends to a continuous section of \(p\) over B.

Let F be a sheaf on B. One says that F is a soft sheaf if the associated étale space (I, p. 50, def. 4) is soft.

Let \(\mathcal{F}\) be a sheaf on B. The sheaf \(\mathcal{F}\) is soft if and only if for every closed set Z of B, every open neighborhood U of Z, and every \(s \in \mathcal{F}(\mathrm{U})\) , there exists \(t \in \mathcal{F}(\mathrm{B})\) and an open neighborhood V of Z contained in U such that \(s \mid \mathrm{V} = t \mid \mathrm{V}\) .

If \(\mathcal{F}\) is a soft sheaf, \(\mathcal{F}(\mathrm{B})\) is nonempty: indeed, the unique section of the étale space \(\mathrm{E}_{\mathcal{F}}\) associated with \(\mathcal{F}\) over \(\varnothing\) extends to a continuous section of \(\mathrm{E}_{\mathcal{F}}\) over B.

Let \(p \colon \mathrm{E} \to \mathrm{B}\) be an étale map and let A be a closed subspace of B. If \(p\) is soft, the map \(p_{\mathrm{A}} \colon \overline{p}^{1}(\mathrm{A}) \to \mathrm{A}\) is soft. Equivalently, if \(\mathcal{F}\) is a soft sheaf on B, the sheaf induced on a closed subspace A is soft.

PROPOSITION 6. --- Let B be a topological space, let \(\mathcal{F}\) be a sheaf on B, and let \((\mathrm{A}_i)_{i\in \mathrm{I}}\) be a locally finite closed covering of B. The sheaf \(\mathcal{F}\) is soft if and only if, for every \(i\in \mathrm{I}\), the induced sheaf \(\mathcal{F}_{\mathrm{A}_i}\) is soft.

The condition is evidently necessary. Let us prove that it is sufficient. Denote by \(p\colon \mathrm{E}\to \mathrm{B}\) the B-étale space \(\mathrm{E}_{\mathcal{F}}\) associated with the sheaf \(\mathcal{F}\) . Let A be a closed subspace of B and \(s\colon \mathrm{A}\to \mathrm{E}\) a continuous section of \(p\) over A; it must be shown that \(s\) has a continuous extension to B. For every subset J of I, set \(\mathrm{A_J} = \bigcup_{i\in \mathrm{J}}\mathrm{A}_i\) ; the set \(\mathrm{A_J}\) is closed in B (TG, I, p. 6, prop. 4).

Let \(\mathcal{S}\) be the set of pairs \((\mathrm{J},t)\) where \(\mathrm{J}\) is a subset of I and \(t\) is a continuous section of E over \(\mathrm{A_J}\) which coincides with \(s\) on \(\mathrm{A} \cap \mathrm{A_J}\) . Equip \(\mathcal{S}\) with the order relation denoted \(\leqslant\) for which \((\mathrm{J},t) \leqslant (\mathrm{J}',t')\) if \(\mathrm{J} \subset \mathrm{J}'\) and \(t' \mid \mathrm{A_J} = t\) . For \(\sigma = (\mathrm{J},t) \in \mathcal{S}\) , write \(\mathrm{J}_{\sigma} = \mathrm{J}\) and \(t_{\sigma} = t\) . Let us show that the ordered set \(\mathcal{S}\) is inductive. Let S be a totally ordered subset of \(\mathcal{S}\) . Set \(\mathrm{J} = \bigcup_{\sigma \in \mathrm{S}} \mathrm{J}_{\sigma}\) ; this is a subset of I. One then defines a section \(t\) of E over \(\mathrm{A_J}\) by setting \(t(x) = t_{\sigma}(x)\) , if \(x \in \mathrm{A}_{J_{\sigma}}\) ; one therefore has \(t \mid \mathrm{A} \cap \mathrm{A_J} = s\) . Let \(j \in \mathrm{J}\) and let \(\sigma \in S\) such that \(j \in J_{\sigma}\) ; since \(t \mid A_j = t_{\sigma} \mid A_j\) the restriction of \(t\) to \(A_j\) is continuous. It then follows from TG, I, p. 19, prop. 4 that \(t\) is continuous. Thus, (J, t) is an element of \(\mathcal{S}\) ; by construction, it is an upper bound of S. This proves that the set \(\mathcal{S}\) is inductive. It therefore possesses a maximal element (J, t) (E, III, p. 20, th. 2).

Let us argue by contradiction, assuming that \(J \neq I\) . Let i be an element of \(I \setminus J\). Set \(A' = (A_i \cap A) \cup (A_i \cap A_J)\) and define a section \(s'\) of E over \(A'\) by:

\[
s ^ {\prime} (a) = \left\{ \begin{array}{l l} s (a) & \text {for} a \in \mathrm{A} _ {i} \cap \mathrm{A}, \\ t (a) & \text {for} a \in \mathrm{A} _ {i} \cap \mathrm{A} _ {\mathrm{J}}, \end{array} \right.
\]

which is possible since \(s\) and \(t\) coincide on \(\mathrm{A} \cap \mathrm{A}_{\mathrm{J}}\) . Moreover, since \(\mathrm{A}_{i} \cap \mathrm{A}\) and \(\mathrm{A}_{i} \cap \mathrm{A}_{\mathrm{J}}\) are closed, the section \(s'\) is continuous (TG, I, p. 19, prop. 4). By hypothesis, there exists a continuous section \(s_{i}: \mathrm{A}_{i} \to \mathrm{E}\) extending \(s'\) . Since the restrictions of \(s_{i}\) and \(t\) to \(\mathrm{A}_{\mathrm{J}} \cap \mathrm{A}_{i}\) are equal, the map \(t': \mathrm{A}_{\mathrm{J} \cup \{i\}} \to \mathrm{E}\) which coincides with \(t\) on \(\mathrm{A}_{\mathrm{J}}\) and with \(s_{i}\) on \(\mathrm{A}_{i}\) is a continuous section of \(p\) over \(\mathrm{A}_{\mathrm{J} \cup \{i\}}\) , extending \(s| \mathrm{A} \cap \mathrm{A}_{\mathrm{J} \cup \{i\}}\) . One then has \((\mathrm{J}, t) < (\mathrm{J} \cup \{i\}, t')\) , which contradicts the hypothesis that \((\mathrm{J}, t)\) is maximal.

Thus, J = I, hence \(A_{J} = B\) and t is a continuous section of E over B extending s.

COROLLARY 1. --- Let B be a paracompact space, \(\mathcal{F}\) a sheaf on B, and \((\mathrm{U}_i)_{i\in \mathrm{I}}\) an open covering of B. If, for every \(i\in \mathrm{I}\) , the induced sheaf \(\mathcal{F}|\mathrm{U}_i\) is soft, then the sheaf \(\mathcal{F}\) is soft.

Indeed, there exists a locally finite closed covering \((\mathrm{F}_{j})_{j\in\mathrm{J}}\) finer than the covering \((\mathrm{U}_{i})_{i\in\mathrm{I}}\) (TG, IX, p. 49, prop. 4 and p. 48, cor. 1). Consequently, for every \(j\in J\), the sheaf \(\mathcal{F}|F_{j}\) is soft and the proposition implies that the sheaf \(\mathcal{F}\) is soft.

COROLLARY 2. --- Let B be a paracompact space, \(\mathcal{F}\) a sheaf on B, and \((\mathrm{A}_i)_{i\in \mathrm{I}}\) a locally finite closed covering of B. For the sheaf \(\mathcal{F}\) to be soft, it is necessary and sufficient that the following condition be satisfied:

For every \(i \in I\) , every closed subset A of \(A_{i}\) , every open set V of B containing A and every element s of \(\mathcal{F}(V)\) , there exists an open neighborhood U of \(A_{i}\) in B, an element t of \(\mathcal{F}(U)\) and an open neighborhood W of A in B contained in \(U \cap V\) such that \(t | W = s | W\) .

Suppose that the sheaf \(\mathcal{F}\) is soft and let us show that the condition is satisfied. Let \(i\in\mathrm{I}\) , let A be a closed subset of \(\mathrm{A}_{i}\) , let V be an open set of B containing A and let \(s\) be an element of \(\mathcal{F}(\mathrm{V})\) . Set \(s_{0}=\sigma_{\mathcal{F}}(s)\) . It is a continuous section over V of the étale space \(\mathrm{E}_{\mathcal{F}}\) . Since \(\mathrm{A}_{i}\) is closed, A is closed in B and \(s_{0}\mid\mathrm{A}\) extends to a section \(t_{0}\colon\mathrm{B}\to\mathrm{E}_{\mathcal{F}}\) , by definition of a soft sheaf. The sections \(s_{0}\) and \(t_{0}\mid\mathrm{V}\) of the V-étale space \(\mathrm{E}_{\mathcal{F}}\times_{\mathrm{B}}\mathrm{V}\) coincide on A, hence on a neighborhood W of A (I, p. 34, prop. 11, b)).

Conversely, suppose that the condition of the corollary is satisfied. By Proposition 6, it suffices to show that, for every \(i \in I\), the sheaf \(\mathcal{F}_{A_i}\) is soft. Let A be a closed subset of \(A_i\) and let \(s_0\) be a section of the étale space of \(\mathcal{F} | A_i\) over A, that is, a continuous section over A of the étale space \(E_{\mathcal{F}}\) . Since A is closed in B and B is paracompact, \(s_0\) extends to a section \(s\) over a neighborhood V of A in B (I, p. 37, th. 2). By hypothesis, there exists an open neighborhood U of \(A_i\) in B and a section \(t\) of \(E_{\mathcal{F}}\) over U which coincides with \(s\) on a neighborhood of A. The restriction of \(t\) to \(A_i\) is a section of \(\mathcal{F}_{A_i}\) which extends \(s_0\) . The sheaf \(\mathcal{F}_{A_i}\) is therefore soft. By Proposition 6, the sheaf \(\mathcal{F}\) is soft.

